\chapter{Complexity — Code Prefix}
%%% generated by the blueprint pipeline from TCSlib.Complexity.TimeHierarchy.CodePrefix
\section{Overview}

Declarations of the Lean module \texttt{TCSlib.Complexity.TimeHierarchy.CodePrefix}.

\section{Declarations}

\begin{definition}[Matching-prefix scan]
\lean{Complexity.TimeHierarchy.scanPre}
\label{Complexity.TimeHierarchy.scanPre}
\leanok
Define a map on Boolean words recursively as follows: it sends the empty word to the
empty word and a one-bit word to itself. On a word beginning with $b,b'$, it retains
these first two bits and then continues recursively on the remaining word if $b=b'$; if
$b\ne b'$, it returns just those first two bits.
\end{definition}

\begin{lemma}[Scanning a paired encoding]
\lean{Complexity.TimeHierarchy.scanPre_pairEncode}
\label{Complexity.TimeHierarchy.scanPre_pairEncode}
\leanok
\uses{Complexity.TimeHierarchy.scanPre, Turing.pairEncode}
\difficulty{4}
For any two finite bit strings $\alpha$ and $w$, applying $\mathrm{scanPre}$ to their pair encoding returns the word obtained by repeating each bit of $\alpha$ twice and then appending the separator $01$.
\end{lemma}

\begin{lemma}[Scanning a nested pair encoding]
\lean{Complexity.TimeHierarchy.scanPre_pairEncode_append}
\label{Complexity.TimeHierarchy.scanPre_pairEncode_append}
\leanok
\uses{Complexity.TimeHierarchy.scanPre, Complexity.TimeHierarchy.scanPre_pairEncode, Turing.pairEncode}
\difficulty{3}
For any two finite bit strings $\alpha$ and $w$, scan the encoding
$\langle\alpha,w\rangle$ from the left in consecutive pairs, retaining each pair until
the first unequal pair, which is retained before the scan stops. Appending
$\langle\alpha,w\rangle$ to the resulting string gives
$\langle\alpha,\langle\alpha,w\rangle\rangle$.
\end{lemma}

\begin{definition}[Preprocessing machine states]
\lean{Complexity.TimeHierarchy.PreState}
\label{Complexity.TimeHierarchy.PreState}
\leanok
Let $PreState$ be the finite state type with states $scan0$, $rwStart$, $rwScan$, and
$copy$, together with a state $scan1(b)$ for each Boolean value $b$.
\end{definition}

\begin{definition}[Preprocessing transition rule]
\lean{Complexity.TimeHierarchy.preTr}
\label{Complexity.TimeHierarchy.preTr}
\leanok
\uses{Complexity.TimeHierarchy.PreState, Turing.Action, Turing.FinTM.controlAction}
For the zero-work-tape Boolean machine, the transition rule is as follows. In state
$scan0$, on reading bit $b$, it moves right, outputs $b$, and enters $scan1(b)$; on a
blank it moves right and enters $rwStart$. In state $scan1(b)$, on reading bit $b'$, it
moves right, outputs $b'$, and enters $scan0$ if $b=b'$, or $rwStart$ otherwise; on a
blank it moves right and enters $rwStart$. In $rwStart$ it moves left into $rwScan$; in
$rwScan$ it moves left on a bit and moves right into $copy$ on a blank. In $copy$ it
moves right and outputs each bit it reads, and halts on a blank.
\end{definition}

\begin{definition}[Preprocessing Turing machine]
\lean{Complexity.TimeHierarchy.preTM}
\label{Complexity.TimeHierarchy.preTM}
\leanok
\uses{Complexity.TimeHierarchy.PreState, Complexity.TimeHierarchy.preTr, Turing.FinTM}
There is a finite Boolean Turing machine with no work tapes, initial state $scan0$, and
transition rule given by the preprocessing transition rule described above.
\end{definition}

\begin{lemma}[Preprocessing machine step]
\lean{Complexity.TimeHierarchy.preTM_step}
\label{Complexity.TimeHierarchy.preTM_step}
\leanok
\uses{Complexity.TimeHierarchy.PreState, Complexity.TimeHierarchy.preTM, Complexity.TimeHierarchy.preTr, Turing.Action.apply, Turing.Cfg, Turing.Cfg.inputSymbol, Turing.Cfg.workTapeSymbols, Turing.MultiTapeTM.step}
\difficulty{3}
Let $x$ be a Boolean string, let $c$ be a configuration of the zero-work-tape
preprocessing machine on input $x$, and let $q$ be a preprocessing state. If the state
of $c$ is $q$, then one step of the machine on $c$ is the result of applying the action
prescribed by the preprocessing transition function at state $q$, with the symbols read
from $c$.
\end{lemma}

\begin{theorem}[Preprocessing scan bound]
\lean{Complexity.TimeHierarchy.pre_scan}
\label{Complexity.TimeHierarchy.pre_scan}
\leanok
\uses{Complexity.TimeHierarchy.PreState, Complexity.TimeHierarchy.preTM, Complexity.TimeHierarchy.preTM_step, Complexity.TimeHierarchy.preTr, Complexity.TimeHierarchy.scanPre, Turing.Action.apply, Turing.Cfg, Turing.Cfg.inputSymbol, Turing.Cfg.workTapeSymbols, Turing.FinTM.controlAction, Turing.FinTM.inputSymbol_at, Turing.MultiTapeTM.runFrom, Turing.MultiTapeTM.runFrom_add, Turing.MultiTapeTM.step, Turing.moveInputPos_pos_of_ne_right}
\difficulty{6}
Let $x$ and $r$ be finite Boolean strings, let $i$ be a natural number, and let $c$ be a
configuration of the zero-work-tape Boolean machine on input $x$. Suppose that removing
the first $i$ bits of $x$ gives $r$, that $i\leq |x|$, that $c$ is in state
$\mathrm{scan0}$, and that its input head is at position $i+1$. For the machine whose
scan state emits each examined bit and, after examining two consecutive unequal bits,
starts its rewrite phase, there is a time $t\leq |r|+1$ when it is in state
$\mathrm{rwStart}$ and its output is the output already in $c$, followed by the string
obtained from $r$ by retaining its first two bits and continuing recursively while each
examined pair is equal (retaining a final singleton if present).
\end{theorem}

\begin{theorem}[Copying the remaining input]
\lean{Complexity.TimeHierarchy.pre_copy}
\label{Complexity.TimeHierarchy.pre_copy}
\leanok
\uses{Complexity.TimeHierarchy.PreState, Complexity.TimeHierarchy.preTM, Complexity.TimeHierarchy.preTM_step, Complexity.TimeHierarchy.preTr, Turing.Action.apply, Turing.Cfg, Turing.Cfg.inputSymbol, Turing.Cfg.workTapeSymbols, Turing.FinTM.inputSymbol_at, Turing.MultiTapeTM.runFrom, Turing.MultiTapeTM.runFrom_succ_eq_step, Turing.MultiTapeTM.step, Turing.moveInputPos_pos_of_ne_right}
\difficulty{5}
Let $x$ and $rest$ be Boolean strings, let $i$ be a natural number, and let $c$ be a
zero-work-tape configuration on input $x$. Suppose that deleting the first $i$ bits of
$x$ gives $rest$, that $i\leq |x|$, that $c$ is in the copy state, and that its input
head is at position $i+1$. After $|rest|+1$ steps, the machine is halted and its output
is the original output of $c$ followed by $rest$.
\end{theorem}

\begin{theorem}[Preprocessing output and time bound]
\lean{Complexity.TimeHierarchy.preTM_computes}
\label{Complexity.TimeHierarchy.preTM_computes}
\leanok
\uses{Complexity.TimeHierarchy.PreState, Complexity.TimeHierarchy.preTM, Complexity.TimeHierarchy.pre_copy, Complexity.TimeHierarchy.pre_scan, Complexity.TimeHierarchy.scanPre, Turing.Cfg, Turing.FinTM.ComputesInTime, Turing.FinTM.ComputesInTime.mono, Turing.FinTM.computesInTime_iff, Turing.FinTM.timed_rewind, Turing.MultiTapeTM.initCfg, Turing.MultiTapeTM.runFrom, Turing.MultiTapeTM.runFrom_add}
\difficulty{5}
For every finite binary word $x$, let $y$ be obtained by retaining the first two bits
and then, in each successive pair, retaining both bits and continuing only when they
agree; if a pair disagrees, discard it and everything after it. The preprocessing
machine outputs the concatenation of $y$ and $x$ within $3|x|+5$ steps.
\end{theorem}
